% [Sumaiya] Data and Exploratory Analysis (rubric: 2.5 marks; the biggest data-science row).
% Top band: source, collection, size, variables described; cleaning steps (missing values, outliers,
%   encoding, leakage) JUSTIFIED; EDA figures that each make a point and VISIBLY inform the modelling.
% FB2 (proposal feedback): show each listed problem being handled: the before -> after table from S2
%   (refused dispatches, zero timestamps, pre-booked, > 3 h durations, missing prep time).
% Numbers: cleaning from notebooks/01_data_audit (all days, before the split); EDA from notebooks/02_eda
%   (training days, 17-22 Oct). README sections 3-4 log every rule (C1-C7) and feature (F1-F9).
\section{Data and Exploratory Analysis}
\label{sec:data}

\subsection{Source and description}
We use the data Meituan released for the INFORMS TSL Data-Driven Research
Challenge~\cite{zhao2024meituan}: every food-delivery waybill logged by its dispatch system in one city
from 17 to 24 October 2022. A row is a \emph{dispatch attempt}, so a refused order appears again when it
is reassigned. The raw file has 654{,}343 rows and 23 columns: 568{,}546 orders, 23 delivery districts
(Meituan's business areas), 4{,}962 restaurants and 4{,}955 couriers. Each row holds the order time,
restaurant and customer coordinates (shifted for privacy, so locations are not real and distances are
approximate), the district, a pre-booking flag, the platform's estimated meal-ready time and promised
delivery time, and the times of dispatch, courier acceptance, pickup and arrival. The target is arrival
minus order time in minutes (median 27.8 on the training days, right-skewed). The file does not say when
the prep estimate is set; we assume it is known at ordering. After-order times and the courier are used
only for analysis (Section~\ref{sec:leakage}); the promise only as a baseline.

\subsection{Cleaning and preprocessing}
Table~\ref{tab:cleaning} lists every rule that removes rows; the reasons follow.

\begin{table}[t]
\centering
\caption{Cleaning keeps 546{,}340 of 654{,}343 rows (all eight days, before the time split); \% is of the
rows before each step. There are no exact duplicate rows. A further 82{,}800 modelling orders (15.2\%)
with a missing (1.4\%) or default (13.8\%) prep estimate are kept, imputed and flagged, not removed.}
\label{tab:cleaning}
\begin{tabular}{lrrr}
\toprule
Step & Removed & \% & Rows left \\
\midrule
Raw dispatch attempts & & & 654{,}343 \\
Keep accepted attempts & 85{,}797 & 13.1 & 568{,}546 \\
Drop orders with no arrival time & 1 & 0.0 & 568{,}545 \\
Exclude pre-booked orders & 22{,}181 & 3.9 & 546{,}364 \\
Exclude deliveries over 180 min & 24 & 0.0 & 546{,}340 \\
\midrule
\multicolumn{4}{l}{Train 17--22 Oct: 406{,}553 \quad Test 23--24 Oct: 139{,}787} \\
\bottomrule
\end{tabular}
\end{table}

\textbf{Refused attempts.} All 85{,}797 repeat an order another courier later accepted, so dropping
them loses no order.

\textbf{Hidden missing values.} The file has no empty cells, but a timestamp of 0 means ``not recorded'':
one order has no arrival time (dropped, since it has no target) and 9{,}874 accepted orders have no prep
estimate (7{,}439 of them in the modelling table). A second missing code hides in the prep column:
13.8\% of the modelling orders have an ``estimate'' less than 5 minutes after the order, and 74\% of these
equal, to the second, the moment the order entered the dispatch system, so they are system defaults, not
forecasts. Flagging them is safe: the default is written within seconds of the order and before any
dispatch (none equals the dispatch time). Both codes are set to missing (82{,}800 modelling orders),
filled with the training median inside each model, and kept as two flags, because these orders are about
2--5 minutes slower than orders with a real estimate at the same distance (shown in the distance figure
below).

\begin{figure}[t]
\centering
\includegraphics[width=\columnwidth]{delivery_time_by_prebook}
\caption{Pre-booked orders follow the slot the customer picked, not the delivery (all days, before
cleaning), and hold almost every delivery over 180 minutes, so they are excluded.}
\label{fig:prebook}
\end{figure}

\textbf{Outliers.} Following the usual outlier triage (error, rare but real, or a different population),
pre-booked orders are a different population (Fig.~\ref{fig:prebook}): their clock runs to a slot the
customer chose (median 80 minutes against 28), and they hold 4{,}073 of the 4{,}097 deliveries over three
hours. Mixing them into training raises the boosting model's validation MAE on normal orders from 6.28
to 6.76 minutes, even keeping only those under three hours, so they are excluded. The IQR and $z$-score
rules would flag 8{,}097 and 3{,}879 ordinary long deliveries on this skewed target, so we keep them; only
24 normal orders take over 180 minutes (99.9\% take under 86), and we remove those as implausible (7 of
them on the test days).

\textbf{Encoding.} Districts are one-hot encoded; restaurant ID is left out because 1{,}229 of its
4{,}962 levels have fewer than 10 orders (see the Methodology).

\textbf{Leakage.} A column is a model input only if it is known when the order is placed. The courier,
the acceptance flag and the dispatch, acceptance and pickup times are all written later, so they are kept
out of every model and used only to measure what they would add (Section~\ref{sec:leakage}). Every step
that learns from data, such as the median used for imputation, is fitted on the training folds alone.

\textbf{What we did not find.} There are no duplicate rows (each accepted order appears once), no
timestamp-order violations in seven checks (placed, dispatched, accepted, picked up, arrived, and no
estimate before the order), no coordinates outside the city's area and no constant columns. These rules
therefore cover every problem the audit found.

\subsection{Exploratory findings that shaped the modelling}
All findings below use the training days only, so the test days stay unseen.

\begin{figure}[t]
\centering
\includegraphics[width=\columnwidth]{delivery_vs_distance_prep}
\caption{Delivery time rises with distance, and orders whose prep estimate is a default (dashed) or
stored as 0 (dotted) run slower at the same distance (training days). This led to the distance feature
and the two missing-estimate flags.}
\label{fig:distance}
\end{figure}

\textbf{Distance is the main driver.} The median delivery time rises about 7.1 minutes per kilometre up
to 3~km, where 93\% of orders fall, and about 5.2 per kilometre beyond (Spearman $\rho = 0.55$;
Fig.~\ref{fig:distance}). The spread grows too: the 20th--80th percentile gap widens from 13 to 24
minutes. We checked for a Simpson reversal by repeating this within each district: the slope stays
positive in all 23, between 4.9 and 7.7 minutes per kilometre. Distance explains about two-thirds of the
variance between district medians, but not all of it: districts 18 and 6 are still 6.2 and 5.5 minutes
faster than the typical order of the same distance. This is why the model uses both distance and
district.

\textbf{Meal peaks are slow.} Delivery time and order volume peak together: at 11:00, with 17\% of
orders, the median is 33.5 minutes against 27.8 overall. The 11:00--11:59 and 17:00--19:59 hours are the
only ones above the overall median, and became the lunch and dinner peak flags.

\textbf{Load matters less than it looks.} Recent district load correlates with delivery time
($\rho = 0.30$) mostly because big districts and meal peaks are busy and slow; within the same district
and hour it is only 0.09. We measure load relative to the district's usual level, so big districts do
not look busy just for being big, and expect it to be a weak driver.

\textbf{The promise is padded.} 35\% of promises sit at a 30-minute minimum, and the median order
arrives 8.3 minutes early, yet 14.7\% arrive late (Fig.~\ref{fig:promise_padding}). The number of
orders arriving late more than halves at exactly eight minutes late (798 orders in the 15~seconds
before, 336 in the 15~seconds after), at every promise level, which suggests an operational rule
applied after ordering. Because the promise is padded, beating it on MAE says little about it as a
promise; the Results also compare promises at the same average length.

\begin{figure}[t]
\centering
\includegraphics[width=\columnwidth]{promise_padding}
\caption{Actual minus promised delivery time (training days): most orders arrive early, and arrivals
drop sharply at +8 minutes, so promises are also compared at equal average length.}
\label{fig:promise_padding}
\end{figure}

\textbf{Two smaller decisions.} The log target removes most of the skew (0.89 to $-0.31$) and the
weekend sits on both sides of the split; the Methodology explains why neither is used.
